% ECONOMICS_PROBLEM: {"id": "D63-396", "title": "PMMS for MMS-feasible pair-demand valuations", "jel_code": "D63", "source_id": "OP-0188"}
% !TeX program = xelatex
\documentclass[11pt,a4paper]{article}
\usepackage[margin=23mm]{geometry}
\usepackage{fontspec}
\setmainfont{Latin Modern Roman}
\usepackage{amsmath,amssymb,mathtools}
\usepackage{xcolor,enumitem,booktabs,longtable,array,xurl,hyperref}
\definecolor{muted}{HTML}{53636C}
\hypersetup{colorlinks=true,urlcolor=blue,linkcolor=blue}
\setlength{\parindent}{0pt}
\setlength{\parskip}{5pt}
\newcommand{\R}{\mathbb R}
\newcommand{\E}{\mathbb E}
\newcommand{\N}{\mathbb N}
\newcommand{\ind}{\mathbf 1}
\newcommand{\Var}{\operatorname{Var}}
\newcommand{\Cov}{\operatorname{Cov}}
\newcommand{\supp}{\operatorname{supp}}
\DeclareMathOperator*{\argmin}{arg\,min}
\DeclareMathOperator*{\argmax}{arg\,max}
\newcommand{\entrymeta}[3]{{\sffamily\small\color{muted}\textbf{\texttt{#1}}\quad|\quad #2\par #3\par}}
\newcommand{\Problemref}[1]{\texttt{#1}}
\newcommand{\JELref}[2]{\texttt{#2} (JEL \texttt{#1})}
\begin{document}
\section*{PMMS for MMS-feasible pair-demand valuations}
\textbf{Problem ID:} \texttt{D63-396}\quad \textbf{JEL:} \texttt{D63}\par
% BEGIN ECONOMICS STATEMENT
\label{problem:OP-0188}
% GAME_THEORY_PROBLEM: {"local_id": "GT-058", "original_number": 58, "kind": "P", "entry_kind": "statement_only", "published": false, "review_required": true, "category": "game_theory"}
\label{game:GT-058}
{\sffamily\small\color{muted}
\textbf{GT-058} | Fair division | \textbf{P}: Open Problem 5, with its valuation assumption corrected.
\par}
\subsection*{Definitions and mathematical statement}
Let $N=[n]$ and $M$ be finite. Each normalized monotone $v_i:2^M\to\mathbb R_{\geq0}$ is pair-demand:
\[
 v_i(S)=\max_{T\subseteq S,\ |T|\leq2}v_i(T).
\]
It is \emph{MMS-feasible} if for every $S\subseteq M$ and every two ordered bipartitions $(X_1,X_2)$ and $(Y_1,Y_2)$ of $S$,
\[
 \max\{v_i(X_1),v_i(X_2)\}\geq
 \min\{v_i(Y_1),v_i(Y_2)\}.
\]
Empty parts are permitted. Define the two-way maximin share on $S$ by
\[
 \mu_i^2(S)=\max_{B\subseteq S}\min\{v_i(B),v_i(S\setminus B)\}.
\]
A complete allocation $A$ is pairwise maximin-share fair (PMMS) if $v_i(A_i)\geq\mu_i^2(A_i\cup A_j)$ for every distinct $i,j$. Determine whether every profile of MMS-feasible pair-demand valuations has a complete PMMS allocation.
\subsection*{Short English statement}
Does pairwise maximin fairness always exist for the broader, non-additive pair-demand class satisfying the source's feasibility property?
\subsection*{Variants and scope boundaries}
MMS-feasible is a property of each valuation on every subset, not the statement that the particular $n$-agent instance admits an MMS allocation. The additive-top-two pair-demand subclass is solved by the source. Partial PMMS permits unallocated goods and is weaker than this target.
\subsection*{Source relationship and status}
This is the source's explicit Open Problem 5. Definition 2.6 is essential: omitting it or replacing it by instance-level MMS existence changes the question. Later nonexistence for other PMMS valuation classes does not automatically settle this subclass.
\subsection*{References}
\begingroup\footnotesize\raggedright
\textbf{[S1]} Jarosław Byrka, Franciszek Malinka, and Tomasz Ponitka (2025). \emph{Probing EFX via PMMS: (Non-)Existence Results in Discrete Fair Division}. arXiv:2507.14957v2. \textit{Verified locator:} Definition 2.6; Section 6; Section 7, Open Problem 5. \url{https://arxiv.org/html/2507.14957v2}
\par\endgroup

\subsection*{Common conventions: Source mathematical conventions}

\textbf{Finite games.} For a finite set $A$, $\Delta(A)$ is its probability
simplex. Players choose independent mixed strategies in Nash equilibrium;
correlation is allowed only when explicitly part of the solution concept.
In a game with payoff functions $u_i$ and strategy profile $x$, an additive
$\varepsilon$ Nash equilibrium satisfies
\[
 u_i(x)\geq u_i(x_i',x_{-i})-\varepsilon
 \quad\text{for every player $i$ and every admissible deviation $x_i'$}.
\]
For numerical approximation, payoffs are normalized as locally specified.
An $\varepsilon$ payoff residual is distinct from distance at most
$\varepsilon$ to the set of exact equilibria. Point convergence, convergence
of empirical play, and convergence in distance to an equilibrium set are
also distinct.

\textbf{Computational representations.} Unless stated otherwise, a complexity
question takes rational data with binary encoding; polynomial time is measured
in total bit length. A fixed-accuracy polynomial algorithm is not necessarily
polynomial in $1/\varepsilon$, and neither is necessarily polynomial in
$\log(1/\varepsilon)$. Succinct games and value, demand, or sampling oracles
must declare their interfaces. Exact real solutions require an explicit
representation; an unspecified list of arbitrary real numbers is not a
finite computational output. A claimed algorithmic barrier is meaningful
only for its specified model and reductions.

\textbf{Dynamic games.} Histories, information, observation, and allowable
randomization are part of the model. A stationary strategy depends only on
the current state; a positional strategy in a deterministic graph game
chooses one action at each controlled vertex. Nash equilibrium at an initial
state and equilibrium after every history are different requirements.
An expectation of a pathwise $\liminf$ payoff, a $\liminf$ of expected averages,
a limiting expected average, and a uniform finite-horizon equilibrium are
different criteria. None is silently substituted for another.

\textbf{Allocation and mechanisms.} A complete allocation assigns every item
exactly once unless otherwise stated. Zero-valued goods, negative values,
capacity constraints, feasibility, lotteries, and copies of goods affect
fairness definitions and are specified locally. Dominant-strategy incentive
compatibility, Bayesian incentive compatibility, universal truthfulness,
and truthfulness in expectation impose different inequalities. Whenever
an objective ratio has zero denominator, its convention is stated or those
instances are excluded.

\textbf{Infinite-dimensional models.} Expectations, integrals, stochastic
equations, and optimization objectives require their local regularity,
integrability, filtrations, and admissible domains. A backward--forward
mean-field system is a boundary-value problem; it is not an initial-value
semiflow without an additional construction. Long-time, large-population,
and vanishing-noise limits require an order of limits and a convergence
topology. If a minimum need not be attained, the objective is its infimum
or supremum and the approximation requirement is stated separately.
% END ECONOMICS STATEMENT
\end{document}
